\begin{figure*}[t]
\centering
\begin{tcolorbox}[title=Cheat Sheet for Disambiguation QA (2/2), colback=gray!10, colframe=black!70, boxrule=0.5pt, top=0pt, bottom=0pt, left=0pt, right=0pt, arc=0.8mm, width=1.0\textwidth]
\scriptsize
\begin{lstlisting}
## 3. **Ambiguity Triggers**
- If both antecedents are equally plausible and the sentence gives no further context, **choose "Ambiguous"**.
- Watch for sentences where both X and Y could have performed the action, received the attribute, or owned the object.

---

## 4. **Special Cues**
- **Gender/Number Agreement:** Make sure the pronoun matches the possible antecedent in gender and number.
- **Role/Profession:** Sometimes, the profession or role makes one antecedent more likely (e.g., only a scientist needs a lab assistant).
- **Typical Scenarios:** Use real-world logic (e.g., a mechanic calls a customer about the customer's car, not their own).

---

## 5. **Quick Reference Table**

| Structure                                      | Most Likely Antecedent | When Ambiguous?                |
|------------------------------------------------|-----------------------|-------------------------------|
| X told Y that [pronoun]...                     | Y (if advice/info)    | If both could be true         |
| X did Y because [pronoun]...                   | X or Y (test both)    | If both make sense            |
| X and Y discuss [pronoun]'s Z                  | Context-dependent     | If both could own Z           |
| X called Y and asked [pronoun] to do Z         | Y                     | If both could be asked        |
| X met with Y at [pronoun]'s office             | Context-dependent     | If both could own office      |
| X did Y because [pronoun] [verb/attribute]     | Context-dependent     | If both fit                   |
| The writer and [pronoun] friends               | The writer            | If only one makes sense       |

---

## 6. **Ambiguity Checklist**
- Both antecedents are grammatically possible.
- Both antecedents are logically possible.
- No context or world knowledge tips the scale.
- If all above are true, **choose "Ambiguous"**.

---

**Use this sheet to reason through each step, especially when both antecedents seem possible!**
\end{lstlisting}
\end{tcolorbox}
\caption{
An example of cheat sheet generated for Disambiguation QA.
This is the second half of the cheat sheet.
}
\label{fig:cheat_sheet_disamb_2}
\end{figure*}